% =========================================================================
\section{Circuito 1: amplificador inversor}\label{sec:inversor}
% =========================================================================
\begin{prelab}[Amplificador inversor]
Usando el modelo de amplificador operacional ideal visto en teoría, encuéntrese la ganancia en
tensión $A_V=\Vout/\Vin$ en función de las resistencias del circuito de la \cref{fig:inversor}. Si
imponemos $R_1=1\,\mathrm{k}\Omega$, determínese el valor de $R_2$ para conseguir $A_V=-10$ a $1\,\mathrm{kHz}$.
\end{prelab}

% Figura: amplificador inversor (circuito 1) con tensiones y corrientes del análisis.
\begin{figure}[ht]
\centering
\begin{circuitikz}[circuito]
  \draw (0,0) node[op amp, ao] (A) {};
  % nudo inversor y red de realimentación
  \draw (A.-) -- ++(-0.6,0) coordinate (N);
  \draw (N) to[R, l_=$R_1$, i<^=$i_1$] ++(-2.4,0) coordinate (E);
  \draw (N) |- ++(0,1.5) coordinate (T);
  \draw (A.out) -- ++(0.7,0) coordinate (O);
  \draw (T) to[R, l=$R_2$, i>^=$i_2$] (T -| O) -- (O);
  % generador
  \draw (E) to[short, -o] ++(-0.3,0) coordinate (Ein) node[above, tens] {$\vin$};
  \draw (Ein) -- ++(-0.9,0) coordinate (G) to[sV, l_=$v_g$] ++(0,-2.0) node[ground] {};
  % entrada no inversora a tierra
  \draw (A.+) -- ++(-0.6,0) coordinate (P) -- ++(0,-0.45) node[ground] {};
  % carga y salida
  \draw (O) to[short, *-o] ++(1.0,0) coordinate (S) node[right, tens] {$\vout$};
  \draw (O) ++(0.5,0) to[R, l=$R_L$] ++(0,-2.0) node[ground] {};
  \fill (N) circle (1.6pt);
  % anotaciones del análisis
  \node[tens, anchor=east, font=\scriptsize] at ($(N)+(-0.08,0.75)$) {$\vm=0$};
  \draw[azulprof, line width=0.3pt] ($(N)+(-0.08,0.72)$) -- ($(N)+(-0.02,0.06)$);
  \node[tens, anchor=east, font=\scriptsize] at ($(P)+(-0.1,-0.25)$) {$\vp=0$};
\end{circuitikz}
\caption{Amplificador inversor (figura~2 del guion), con $R_1=1\,\mathrm{k}\Omega$,
$R_L=10\,\mathrm{k}\Omega$ y la alimentación omitida. Se indican las magnitudes del análisis: el nudo
inversor es una tierra virtual ($\vm=\vp=0$) y no absorbe corriente, de modo que $i_1=i_2$.}
\label{fig:inversor}
\end{figure}


\begin{resolucion}
AO ideal con realimentación negativa a través de $R_2$; sentidos de corriente de la \cref{fig:inversor}.

\textbf{1. Tierra virtual.} La entrada $+$ está a tierra; por (R2),
\begin{equation}\label{eq:i-vm}
  \vm=\vp=0 .
\end{equation}

\textbf{2. Ley de Ohm en $R_1$}, entre $\vin$ y $\vm$:
\begin{equation}\label{eq:i-i1}
  i_1=\frac{\vin-\vm}{R_1}=\frac{\vin}{R_1} .
\end{equation}

\textbf{3. Kirchhoff en el nudo inversor}, con $i_-=0$ por (R1):
\begin{equation}\label{eq:i-i2}
  i_1=i_2+i_-\quad\Longrightarrow\quad i_2=i_1=\frac{\vin}{R_1} .
\end{equation}

\textbf{4. Ley de Ohm en $R_2$}, recorrida por $i_2$ del nudo inversor a la salida:
\begin{equation}\label{eq:i-vout}
  \vm-\vout=R_2\,i_2\quad\Longrightarrow\quad \vout=-\frac{R_2}{R_1}\,\vin .
\end{equation}

\textbf{Resultado.}
\begin{equation}\label{eq:i-Av}
  \resultado{A_V=\frac{\vout}{\vin}=-\frac{R_2}{R_1}}
\end{equation}

\textbf{Diseño.} Con $A_V=-10$ y $R_1=1\,\mathrm{k}\Omega$ en \eqref{eq:i-Av}:
\begin{equation}\label{eq:i-R2}
  R_2=-A_V\,R_1=10\cdot1\,\mathrm{k}\Omega\quad\Longrightarrow\quad\resultado{R_2=10\,\mathrm{k}\Omega}
\end{equation}

El signo menos indica contrafase ($180^{\circ}$). La ganancia no depende de la frecuencia, de modo que el
diseño vale a $1\,\mathrm{kHz}$, ni de $A$ ni de $R_L$: solo del cociente $R_2/R_1$.
\end{resolucion}
